\documentclass[11pt,a4paper]{article}

\usepackage{amsmath,amssymb,amsthm}
\usepackage{mathtools}
\usepackage{bm}
\usepackage{enumerate}
\usepackage{hyperref}
\usepackage{cleveref}
\usepackage[margin=1in]{geometry}
\usepackage{tcolorbox}

% Theorem environments
\newtheorem{theorem}{Theorem}[section]
\newtheorem{lemma}[theorem]{Lemma}
\newtheorem{proposition}[theorem]{Proposition}
\newtheorem{corollary}[theorem]{Corollary}
\theoremstyle{definition}
\newtheorem{definition}[theorem]{Definition}
\newtheorem{example}[theorem]{Example}
\newtheorem{remark}[theorem]{Remark}

% Notation shortcuts
\newcommand{\E}{\mathbb{E}}
\newcommand{\R}{\mathbb{R}}
\newcommand{\Var}{\mathrm{Var}}
\newcommand{\Normal}{\mathcal{N}}
\newcommand{\given}{\,|\,}
\newcommand{\btheta}{\boldsymbol{\theta}}
\newcommand{\bxi}{\boldsymbol{\xi}}
\newcommand{\bpsi}{\boldsymbol{\psi}}
\DeclareMathOperator{\CE}{CE}
\DeclareMathOperator{\RDU}{RDU}
\DeclareMathOperator{\diag}{diag}
\DeclareMathOperator{\sgn}{sgn}
\DeclareMathOperator{\tr}{tr}

% Highlight box for key results
\newtcolorbox{keyresult}[1][]{
    colback=blue!5!white,
    colframe=blue!50!black,
    fonttitle=\bfseries,
    title={#1},
}

\title{The Gale--Nikaido Univalence Theorem\\
Applied to RDU Certainty Equivalent Mappings}
\author{Technical Note}
\date{February 2026}

\begin{document}

\maketitle

\begin{abstract}
We apply the Gale--Nikaido univalence theorem to the certainty equivalent
(CE) mapping arising from a general Rank-Dependent Utility (RDU) model.
The mapping takes preference parameters---utility parameters
$\btheta \in \R^{p_1}$ and probability weighting function (PWF) parameters
$\bpsi \in \R^{p_2}$---to a vector of certainty equivalents
$(\CE_1, \ldots, \CE_K)$ for $K = p_1 + p_2$ reference lotteries.
We derive the Jacobian of this mapping in closed form for arbitrary
utility $u(x;\btheta)$ and PWF $w(p;\bpsi)$, identify the P-matrix
condition as the operative constraint for global injectivity, and show
that \emph{prize diversity} across reference lotteries is the economic
mechanism that generically delivers the P-matrix property.
This corrects the companion paper's current statement that the
Gale--Nikaido theorem ``does not directly apply'' to CE mappings:
the theorem's domain requirement is a rectangular region (not
properness, which is Hadamard's condition), and economic parameter
spaces are naturally rectangular.
\end{abstract}

\tableofcontents
\newpage


%==============================================================================
\section{The Gale--Nikaido Univalence Theorem}
\label{sec:gn_theorem}
%==============================================================================

We begin with a precise statement of the theorem and the definitions it requires.

\begin{definition}[P-matrix]
\label{def:p_matrix}
A real $n \times n$ matrix $A$ is a \textbf{P-matrix} if every principal minor
of $A$ is strictly positive. That is, for every nonempty index set
$I \subseteq \{1, \ldots, n\}$, the submatrix $A_{II}$ (obtained by selecting
rows and columns indexed by $I$) satisfies $\det(A_{II}) > 0$.
\end{definition}

For an $n \times n$ matrix, the P-matrix condition requires $2^n - 1$ strict
inequalities: one for each nonempty subset of $\{1, \ldots, n\}$. The conditions
are nested: all $1 \times 1$ principal minors (the diagonal entries) must be
positive, all $\binom{n}{2}$ of the $2 \times 2$ principal minors must be
positive, and so on up to the full determinant.

\begin{remark}[Equivalent characterizations]
\label{rem:p_matrix_equiv}
A matrix $A$ is a P-matrix if and only if, for every nonzero vector
$x \in \R^n$, there exists an index $i$ such that $x_i (Ax)_i > 0$.
This ``sign-reversal'' characterization~\cite{gale_nikaido_1965} is often
more intuitive: no matter what direction you perturb the input, at least one
output component moves in the \emph{same} direction as its corresponding
input component.
\end{remark}

The central result is:

\begin{theorem}[Gale--Nikaido, 1965]
\label{thm:gale_nikaido}
Let $\Omega \subseteq \R^n$ be an open \textbf{rectangular} region (i.e., a
Cartesian product of open intervals: $\Omega = \prod_{i=1}^n (a_i, b_i)$
where $-\infty \leq a_i < b_i \leq +\infty$). Let $\mathbf{g}: \Omega \to \R^n$
be a $C^1$ mapping. If the Jacobian matrix $J_{\mathbf{g}}(\bxi)$ is a
\textbf{P-matrix} for all $\bxi \in \Omega$, then $\mathbf{g}$ is
\textbf{globally injective} (univalent) on $\Omega$.
\end{theorem}

Several features of this theorem deserve emphasis.

\paragraph{Domain requirement: rectangular, not convex.}
The original Gale--Nikaido theorem~\cite{gale_nikaido_1965} states the result
for rectangular domains, which is a special case of convexity but carries an
important distinction. A rectangular domain (product of intervals) is the
natural shape of an economic parameter space: each parameter has its own range,
independent of the others. The domain need not be bounded---each interval
$(a_i, b_i)$ may extend to $\pm\infty$.

\paragraph{No properness required.}
The Gale--Nikaido theorem does \emph{not} require the mapping $\mathbf{g}$ to be
proper (i.e., $\|\mathbf{g}(\bxi)\| \to \infty$ as $\|\bxi\| \to \infty$).
This is a crucial distinction from Hadamard's global inverse function theorem,
which does require properness. CE mappings are bounded (by the prize range)
while parameter domains are typically unbounded, so properness fails
generically for CE mappings. The Gale--Nikaido theorem bypasses this
obstruction entirely by imposing a stronger algebraic condition (P-matrix) on
the Jacobian instead.

\paragraph{P-matrix vs.\ nonsingular Jacobian.}
A P-matrix is nonsingular (since $\det A > 0$ is one of the conditions), but
the converse is false. The P-matrix condition is strictly stronger than
requiring $\det J \neq 0$ everywhere. Whereas a nonsingular Jacobian gives
only local injectivity (by the inverse function theorem), the P-matrix condition
bridges the gap to global injectivity on rectangular domains.

\paragraph{The Gale--Nikaido--Inada refinement.}
Inada~\cite{gale_nikaido_1965} showed that if $\Omega$ is a rectangular domain
and $J_{\mathbf{g}}$ has all diagonal entries positive and is a \emph{weakly
sign-symmetric} P-matrix, then weaker conditions on the off-diagonal entries
suffice. We do not use this refinement here, but it is available for models
where the full P-matrix condition is difficult to verify while sign structure
is transparent.


%==============================================================================
\section{General RDU Model and the CE Mapping}
\label{sec:model}
%==============================================================================

\subsection{Primitives}

Consider a general RDU model with:
\begin{itemize}
    \item \textbf{Utility function:} $u(x; \btheta)$, where $x \geq 0$ is a
    monetary outcome and $\btheta = (\theta_1, \ldots, \theta_{p_1}) \in
    \Theta \subseteq \R^{p_1}$ are utility parameters. We assume $u$ is
    strictly increasing in $x$ for all $\btheta \in \Theta$, i.e.,
    $\partial u / \partial x > 0$.

    \item \textbf{Probability weighting function (PWF):}
    $w(p; \bpsi)$, where $p \in [0,1]$ is a cumulative probability and
    $\bpsi = (\psi_1, \ldots, \psi_{p_2}) \in \Psi \subseteq \R^{p_2}$ are
    PWF parameters. We assume $w(0;\bpsi) = 0$, $w(1;\bpsi) = 1$, and $w$ is
    strictly increasing in $p$ for all $\bpsi \in \Psi$.
\end{itemize}

The combined parameter vector is $\bxi = (\btheta, \bpsi) \in
\Xi = \Theta \times \Psi \subseteq \R^p$ where $p = p_1 + p_2$.

\paragraph{Monotone influence assumptions.}
We require:
\begin{enumerate}
    \item For each utility parameter $\theta_i$, $i = 1, \ldots, p_1$:
    $\partial u(x; \btheta) / \partial \theta_i \neq 0$ for all
    $x > 0$ and $\btheta \in \Theta$. That is, each $\theta_i$ has a
    nontrivial effect on utility at every outcome level.

    \item For each PWF parameter $\psi_j$, $j = 1, \ldots, p_2$:
    $\partial w(p; \bpsi) / \partial \psi_j \neq 0$ for all
    $p \in (0,1)$ and $\bpsi \in \Psi$. That is, each $\psi_j$ has a
    nontrivial effect on the weighting function at every probability level.
\end{enumerate}

These are natural requirements: they simply say that each parameter
\emph{does something}. A parameter that has zero marginal effect on utility (or
weighting) everywhere is redundant and should be dropped from the model.

\subsection{RDU Valuation}

Consider $K$ reference lotteries $L_1, \ldots, L_K$. Each lottery $L_k$ has
$M_k$ outcomes with prizes $x_{1k} < x_{2k} < \cdots < x_{M_k,k}$ and
cumulative probabilities. Under RDU, the value of lottery $L_k$ is:
\begin{equation}
    \label{eq:rdu}
    \RDU_k(\bxi) = \sum_{m=1}^{M_k} \pi_{mk}(\bpsi) \, u(x_{mk}; \btheta)
\end{equation}
where the decision weights $\pi_{mk}(\bpsi)$ are derived from the PWF applied
to cumulative probabilities. For a lottery with outcomes ranked
$x_{1k} < \cdots < x_{M_k,k}$:
\begin{equation}
    \pi_{mk}(\bpsi) = w\!\left(\sum_{j=m}^{M_k} p_{jk};\, \bpsi\right)
    - w\!\left(\sum_{j=m+1}^{M_k} p_{jk};\, \bpsi\right)
\end{equation}
where $p_{jk}$ is the probability of outcome $x_{jk}$ in lottery $L_k$, and
$w(\cdot;\bpsi)$ is the PWF. The weights satisfy $\pi_{mk} > 0$ (by strict
monotonicity of $w$) and $\sum_m \pi_{mk} = 1$.

\subsection{The CE Mapping}

The certainty equivalent $\CE_k$ of lottery $L_k$ is defined implicitly by:
\begin{equation}
    \label{eq:ce_implicit}
    u(\CE_k; \btheta) = \RDU_k(\bxi)
\end{equation}
Since $u$ is strictly increasing in $x$, for each $\bxi$ there is a unique
$\CE_k \in [\min_m x_{mk}, \max_m x_{mk}]$ satisfying~\eqref{eq:ce_implicit}.
We can write $\CE_k = u^{-1}(\RDU_k(\bxi); \btheta)$, where $u^{-1}$ denotes
the inverse of $u$ with respect to its first argument.

The \textbf{CE mapping} is the function:
\begin{equation}
    \label{eq:g_map}
    g: \Xi \to \R^K, \qquad g(\bxi) = (\CE_1(\bxi), \ldots, \CE_K(\bxi))
\end{equation}

For the Jacobian correction methodology to be valid, we need $g$ to be
injective on $\Xi$. A necessary condition is $K \geq p = p_1 + p_2$ (at
least as many lotteries as parameters). In the square case $K = p$, we apply
the Gale--Nikaido theorem.


%==============================================================================
\section{Jacobian of the CE Mapping}
\label{sec:jacobian}
%==============================================================================

Our goal in this section is to compute how the certainty equivalents
$(\CE_1, \ldots, \CE_K)$ change when we perturb the model parameters $\bxi =
(\btheta, \bpsi)$. These rates of change are the entries of the Jacobian
matrix $J_g$. We derive them by implicit differentiation of the defining
equation $u(\CE_k; \btheta) = \RDU_k(\btheta, \bpsi)$, and then organize
the result into a matrix factorization that will be convenient for the
P-matrix analysis in Section~\ref{sec:p_matrix}.

\subsection{Implicit Differentiation}

Recall the defining equation for each CE (equation~\eqref{eq:ce_implicit}):
\begin{equation}
    u(\CE_k;\, \btheta) = \RDU_k(\btheta, \bpsi)
    = \sum_{m=1}^{M_k} \pi_{mk}(\bpsi) \, u(x_{mk};\, \btheta)
    \tag{\ref{eq:ce_implicit}$'$}
\end{equation}
This says: the utility of receiving $\CE_k$ for sure equals the
decision-weighted average utility of lottery $L_k$'s outcomes.

We need to differentiate both sides with respect to each parameter
$\xi_\ell$. The two types of parameters---utility parameters $\theta_i$ and
PWF parameters $\psi_j$---enter the equation differently, so we handle them
separately.

\subsubsection{Derivatives with respect to utility parameters $\theta_i$}

Consider perturbing a utility parameter $\theta_i$ by a small amount. Both
sides of the CE equation change, but for different reasons:
\begin{itemize}
    \item \textbf{Left-hand side} $u(\CE_k;\, \btheta)$: This changes in
    \emph{two} ways. First, the utility function's shape changes (because
    $\theta_i$ enters $u$ directly). Second, $\CE_k$ itself shifts (because
    the RDU value changes, and the CE must adjust to keep the equation
    balanced). By the chain rule:
    \begin{equation}
        \frac{d}{d\theta_i}\!\Big[u(\CE_k;\, \btheta)\Big]
        = \underbrace{
            \frac{\partial u}{\partial x}\bigg|_{x=\CE_k}
            \cdot \frac{\partial \CE_k}{\partial \theta_i}
        }_{\substack{\text{indirect effect:}\\\text{CE shifts}}}
        \;+\;
        \underbrace{
            \frac{\partial u}{\partial \theta_i}\bigg|_{x=\CE_k}
        }_{\substack{\text{direct effect:}\\\text{utility shape changes}}}
    \end{equation}

    \item \textbf{Right-hand side}
    $\sum_m \pi_{mk}(\bpsi)\, u(x_{mk};\, \btheta)$: The decision weights
    $\pi_{mk}$ depend only on $\bpsi$, not on $\btheta$, so they are
    constants with respect to $\theta_i$. The prizes $x_{mk}$ are fixed data.
    Only $u$ changes:
    \begin{equation}
        \frac{d}{d\theta_i}\!\left[\sum_m \pi_{mk}\, u(x_{mk};\, \btheta)
        \right]
        = \sum_{m=1}^{M_k} \pi_{mk}(\bpsi) \,
        \frac{\partial u}{\partial \theta_i}(x_{mk};\, \btheta)
    \end{equation}
\end{itemize}

Setting the two sides equal:
\begin{equation}
    \frac{\partial u}{\partial x}\bigg|_{x=\CE_k} \cdot
    \frac{\partial \CE_k}{\partial \theta_i}
    + \frac{\partial u}{\partial \theta_i}\bigg|_{x=\CE_k}
    = \sum_{m=1}^{M_k} \pi_{mk}(\bpsi) \,
    \frac{\partial u}{\partial \theta_i}(x_{mk}; \btheta)
\end{equation}

Solving for $\partial \CE_k / \partial \theta_i$ (move the direct effect to
the right, divide by $\partial u / \partial x$):
\begin{equation}
    \label{eq:dCE_dtheta}
    \frac{\partial \CE_k}{\partial \theta_i}
    = \frac{
        \overbrace{
            \displaystyle\sum_{m=1}^{M_k} \pi_{mk}(\bpsi) \,
            \frac{\partial u}{\partial \theta_i}(x_{mk}; \btheta)
        }^{\substack{\text{weighted avg of } \partial u/\partial\theta_i\\
        \text{across lottery outcomes}}}
        \;-\;
        \overbrace{
            \frac{\partial u}{\partial \theta_i}(\CE_k; \btheta)
        }^{\substack{\partial u/\partial\theta_i\\
        \text{at the CE}}}
    }{
        \underbrace{
            \dfrac{\partial u}{\partial x}(\CE_k; \btheta)
        }_{> 0 \text{ (u increasing in } x\text{)}}
    }
\end{equation}

The denominator is always positive (utility is strictly increasing in money),
so the sign and magnitude of $\partial \CE_k / \partial \theta_i$ are
determined entirely by the numerator.

\paragraph{Interpreting the numerator.}
The numerator compares two quantities:
\begin{enumerate}
    \item The decision-weighted average of ``how much $\theta_i$ affects
    utility'' across the lottery outcomes $x_{1k}, \ldots, x_{M_k,k}$.
    \item The same quantity evaluated at the single point $x = \CE_k$.
\end{enumerate}

If the function $x \mapsto \partial u / \partial \theta_i(x;\, \btheta)$ is
\emph{linear} in $x$, then the weighted average over the lottery prizes
equals the value at the CE (by definition of the CE as the certainty
equivalent under $u$), and the numerator is zero---the CE does not respond to
$\theta_i$.

But if $\partial u / \partial \theta_i$ is \emph{nonlinear} in $x$ (which is
the generic case for any interesting utility function), then a Jensen's
inequality gap opens: the weighted average over prizes differs from the value
at the CE, and the CE must shift in response to a change in $\theta_i$.
This Jensen's inequality mechanism is what makes the utility parameter
derivatives generically nonzero.

\subsubsection{Derivatives with respect to PWF parameters $\psi_j$}

Now consider perturbing a PWF parameter $\psi_j$. The story is simpler
because $\psi_j$ does not affect the utility function directly---it only
changes the decision weights $\pi_{mk}(\bpsi)$.

\begin{itemize}
    \item \textbf{Left-hand side} $u(\CE_k;\, \btheta)$: The utility function
    $u$ does not depend on $\bpsi$, so the only change comes through the
    CE shifting:
    \begin{equation}
        \frac{d}{d\psi_j}\!\Big[u(\CE_k;\, \btheta)\Big]
        = \frac{\partial u}{\partial x}\bigg|_{x=\CE_k}
        \cdot \frac{\partial \CE_k}{\partial \psi_j}
    \end{equation}
    There is no ``direct effect'' term because $u$ itself is unaffected by
    $\psi_j$.

    \item \textbf{Right-hand side}
    $\sum_m \pi_{mk}(\bpsi)\, u(x_{mk};\, \btheta)$: Now the utility values
    $u(x_{mk};\, \btheta)$ are fixed (they don't depend on $\bpsi$), but the
    weights $\pi_{mk}(\bpsi)$ change:
    \begin{equation}
        \frac{d}{d\psi_j}\!\left[\sum_m \pi_{mk}\, u(x_{mk})
        \right]
        = \sum_{m=1}^{M_k}
        \frac{\partial \pi_{mk}}{\partial \psi_j}(\bpsi) \, u(x_{mk};\, \btheta)
    \end{equation}
\end{itemize}

Setting equal and solving:
\begin{equation}
    \label{eq:dCE_dpsi}
    \frac{\partial \CE_k}{\partial \psi_j}
    = \frac{
        \displaystyle\sum_{m=1}^{M_k}
        \frac{\partial \pi_{mk}}{\partial \psi_j}(\bpsi) \,
        u(x_{mk}; \btheta)
    }{
        \dfrac{\partial u}{\partial x}(\CE_k; \btheta)
    }
\end{equation}

\paragraph{Interpreting the numerator.}
The key observation is that the decision weights always sum to one:
$\sum_m \pi_{mk} = 1$ for all $\bpsi$. Differentiating this constraint gives
$\sum_m \partial\pi_{mk}/\partial\psi_j = 0$. So the weight perturbations
$(\partial\pi_{1k}/\partial\psi_j, \ldots,
\partial\pi_{M_k,k}/\partial\psi_j)$ sum to zero: when $\psi_j$ increases
weight on some outcomes, it must decrease weight on others by the same total
amount.

The numerator of~\eqref{eq:dCE_dpsi} is the inner product of these
mean-zero weight perturbations with the utility values. In plain language:
\emph{does $\psi_j$ shift weight toward high-utility outcomes or low-utility
outcomes?} If $\psi_j$ shifts weight toward outcomes with above-average
utility, the RDU value increases, and the CE must increase to match. If
$\psi_j$ shifts weight toward low-utility outcomes, the CE decreases.

The numerator is zero only if the weight perturbation is ``orthogonal'' to
the utility vector---meaning $\psi_j$ reshuffles weight among outcomes with
equal utility, which is non-generic.


\subsection{Matrix Factorization: Separating the Common Denominator}

Looking at equations~\eqref{eq:dCE_dtheta} and~\eqref{eq:dCE_dpsi}, every
entry of the Jacobian $J_g$ (i.e., every $\partial \CE_k / \partial \xi_\ell$)
has the same denominator: $\partial u / \partial x$ evaluated at $\CE_k$.
This denominator depends on $k$ (different lotteries have different CEs) but
not on $\ell$ (which parameter we differentiate with respect to). We can
therefore factor it out row by row.

Define the positive scalar for each lottery $k$:
\begin{equation}
    d_k(\bxi) \coloneqq \frac{\partial u}{\partial x}(\CE_k; \btheta) > 0
\end{equation}
This is positive because $u$ is strictly increasing in $x$---a dollar more
always increases utility. Collect these into a diagonal matrix
$D = \diag(d_1, \ldots, d_K)$.

Every entry of the Jacobian can be written as
$(\partial \CE_k / \partial \xi_\ell) = (\text{numerator}_{k\ell}) / d_k$.
In matrix form:
\begin{equation}
    \label{eq:factorization}
    J_g(\bxi) = D(\bxi)^{-1} \, \widetilde{J}(\bxi)
\end{equation}
where $\widetilde{J}$ is the ``unnormalized'' Jacobian whose entries are just
the numerators:
\begin{equation}
    \widetilde{J}_{k,\ell} = d_k \cdot \frac{\partial \CE_k}{\partial \xi_\ell}
    \qquad\text{(numerator of the $\partial\CE_k / \partial\xi_\ell$ formula)}
\end{equation}

The point of this factorization is that $D$ is a positive diagonal matrix,
and---as we will show in Section~\ref{sec:p_matrix}---multiplying by a
positive diagonal matrix does not affect the P-matrix property. This means
we can forget about the denominators $d_k$ entirely and work with
$\widetilde{J}$, which has simpler entries.

In the square case $K = p$, we write $\widetilde{J}$ in block form
corresponding to the two types of parameters:
\begin{equation}
    \label{eq:J_tilde_block}
    \widetilde{J} = \begin{pmatrix} A & B \end{pmatrix}
\end{equation}
where $A \in \R^{K \times p_1}$ contains the numerators from the utility
parameter derivatives, and $B \in \R^{K \times p_2}$ contains the numerators
from the PWF parameter derivatives. Each row of $\widetilde{J}$ corresponds to
one lottery; the first $p_1$ columns correspond to utility parameters and the
last $p_2$ columns to PWF parameters.

\subsection{The Utility Parameter Block $A$}

From~\eqref{eq:dCE_dtheta}, the $(k, i)$ entry of $A$ is just the numerator
of $\partial \CE_k / \partial \theta_i$:
\begin{equation}
    \label{eq:A_entry}
    A_{ki} = \sum_{m=1}^{M_k} \pi_{mk}(\bpsi) \,
    \frac{\partial u}{\partial \theta_i}(x_{mk}; \btheta)
    - \frac{\partial u}{\partial \theta_i}(\CE_k; \btheta)
\end{equation}

Reading this in words: $A_{ki}$ is the \emph{gap} between the
decision-weighted average of $\partial u / \partial \theta_i$ across lottery
$L_k$'s prizes and the same derivative evaluated at the certainty equivalent.
This is the Jensen's-inequality gap discussed above.

\begin{remark}[When is $A_{ki}$ nonzero?]
\label{rem:A_interp}
$A_{ki} = 0$ would require a very specific cancellation: the weighted average
of $\partial u / \partial \theta_i$ over the lottery prizes would need to
\emph{exactly equal} the value at the CE. For a utility function where
$\partial u / \partial \theta_i$ is nonlinear in $x$, this is a non-generic
(knife-edge) condition. So generically, $A_{ki} \neq 0$ for all
lottery-parameter pairs.

As a concrete example: for CRRA utility $u(x; r) = x^{1-r}/(1-r)$, the
derivative $\partial u / \partial r$ involves terms like $x^{1-r} \ln x$,
which is nonlinear in $x$. The weighted average over prizes therefore
differs from the value at the CE.
\end{remark}

\subsection{The PWF Parameter Block $B$}

From~\eqref{eq:dCE_dpsi}, the $(k, j)$ entry of $B$ is just the numerator
of $\partial \CE_k / \partial \psi_j$:
\begin{equation}
    \label{eq:B_entry}
    B_{kj} = \sum_{m=1}^{M_k}
    \frac{\partial \pi_{mk}}{\partial \psi_j}(\bpsi) \,
    u(x_{mk}; \btheta)
\end{equation}

Reading this in words: $B_{kj}$ measures how the PWF parameter $\psi_j$
\emph{reshuffles} the decision weights of lottery $L_k$'s outcomes, weighted
by the utility of each outcome.

\begin{remark}[When is $B_{kj}$ nonzero?]
\label{rem:B_interp}
Since the weight perturbation sums to zero
($\sum_m \partial\pi_{mk}/\partial\psi_j = 0$), $B_{kj}$ is zero only if
the perturbation is ``orthogonal'' to the utility vector---that is, if
$\psi_j$ shifts weight between outcomes that happen to have exactly the same
utility. With distinct prizes and a strictly monotone $u$, the utility values
$u(x_{1k}) < u(x_{2k}) < \cdots < u(x_{M_k,k})$ are all different, so any
nonzero weight reshuffling produces a nonzero $B_{kj}$.

As a concrete example: for the Prelec PWF, increasing $\eta$ shifts weight
away from high-probability outcomes toward low-probability outcomes (it
compresses the weighting function toward zero). Whether this increases or
decreases $B_{kj}$ depends on whether the low-probability outcomes are the
high-utility or low-utility ones---which depends on the specific lottery
design.
\end{remark}


%==============================================================================
\section{The P-Matrix Condition}
\label{sec:p_matrix}
%==============================================================================

With the Jacobian's structure in hand, we now connect it to the
Gale--Nikaido theorem. The theorem requires the Jacobian to be a P-matrix
at every point. This section shows that the factorization from
Section~\ref{sec:jacobian} lets us ignore the denominators $d_k$ and work
directly with the simpler matrix $\widetilde{J}$, then unpacks what the
P-matrix condition actually requires of the entries.

\subsection{Reduction to $\widetilde{J}$: the Denominators Do Not Matter}

Recall that $J_g = D^{-1} \widetilde{J}$, where $D = \diag(d_1, \ldots, d_K)$
with each $d_k > 0$. The following lemma says that this positive diagonal
scaling is irrelevant for the P-matrix property.

\begin{lemma}
\label{lem:D_reduction}
$J_g(\bxi)$ is a P-matrix if and only if $\widetilde{J}(\bxi)$ is a P-matrix.
\end{lemma}

\begin{proof}
We need to show that multiplying on the left by $D^{-1}$ (a positive diagonal
matrix) preserves the sign of every principal minor. Pick any nonempty index
set $I \subseteq \{1, \ldots, K\}$ and consider the principal submatrix
$(D^{-1}\widetilde{J})_{II}$---the submatrix obtained by selecting rows and
columns indexed by $I$. Because $D^{-1}$ is diagonal, the rows of
$D^{-1}\widetilde{J}$ are just the rows of $\widetilde{J}$ each scaled by
$d_k^{-1}$. So the submatrix decomposes as:
\begin{equation}
    (D^{-1}\widetilde{J})_{II} = D_I^{-1} \cdot \widetilde{J}_{II}
\end{equation}
where $D_I^{-1} = \diag(d_k^{-1} : k \in I)$. Taking determinants:
\begin{equation}
    \det\!\big((D^{-1}\widetilde{J})_{II}\big)
    = \underbrace{\prod_{k \in I} d_k^{-1}}_{> 0}
    \;\cdot\; \det(\widetilde{J}_{II})
\end{equation}
The prefactor is a product of positive numbers, hence positive. So
$\det((D^{-1}\widetilde{J})_{II}) > 0$ if and only if
$\det(\widetilde{J}_{II}) > 0$. Since this holds for \emph{every} index set
$I$, the two matrices have identical principal minor signs, and one is a
P-matrix if and only if the other is.
\end{proof}

The practical upshot: we can forget about the denominators
$d_k = \partial u / \partial x |_{\CE_k}$ and analyze the P-matrix property
using $\widetilde{J} = [A \mid B]$ alone, which has the simpler entries
derived in Section~\ref{sec:jacobian}.

\subsection{What the P-Matrix Condition Requires: Unpacking the $2^p - 1$
Conditions}

For $K = p$ lotteries and $p$ parameters, the P-matrix condition on
$\widetilde{J}$ requires all $2^p - 1$ principal minors to be strictly
positive. To understand what this means concretely, it helps to think about
the conditions layer by layer, from the simplest (order 1) to the most
complex (order $p$).

\paragraph{Layer 1: Diagonal entries ($p$ conditions).}
The order-1 principal minors are just the diagonal entries:
$\widetilde{J}_{kk} > 0$ for $k = 1, \ldots, p$. In the block
structure $\widetilde{J} = [A \mid B]$, the diagonal entries come from both
blocks depending on how lotteries are paired with parameters.

To make this concrete, think of it as an ``assignment'': lottery $k$ is
assigned to parameter $\xi_k$. The diagonal condition says each lottery must
respond \emph{positively} to its assigned parameter---when $\xi_k$ increases,
$\CE_k$ must increase (after accounting for the common denominator $d_k$).

If even one diagonal entry is zero or negative, the P-matrix condition fails
immediately.

\paragraph{Layer 2: $2 \times 2$ principal minors ($\binom{p}{2}$ conditions).}
For each pair of indices $(k, \ell)$, extract the $2 \times 2$ submatrix
from rows and columns $k, \ell$:
\begin{equation}
    \begin{pmatrix}
        \widetilde{J}_{kk} & \widetilde{J}_{k\ell} \\
        \widetilde{J}_{\ell k} & \widetilde{J}_{\ell\ell}
    \end{pmatrix}
\end{equation}
Its determinant must be positive:
\begin{equation}
    \widetilde{J}_{kk}\widetilde{J}_{\ell\ell} -
    \widetilde{J}_{k\ell}\widetilde{J}_{\ell k} > 0
\end{equation}

In words: the product of the two ``own effects'' (lottery $k$'s response to
parameter $k$, times lottery $\ell$'s response to parameter $\ell$) must
exceed the product of the two ``cross effects'' (lottery $k$'s response to
parameter $\ell$, times lottery $\ell$'s response to parameter $k$).

This is a ``not too much cross-talk'' condition. It allows cross-effects to
exist---changing a PWF parameter will generally affect all lotteries' CEs,
not just the one ``assigned'' to it---but the cross-effects cannot be so
large that they overwhelm the own-effects. The own-effects must dominate in a
pairwise sense.

\paragraph{Layer 3 and beyond: Higher-order principal minors.}
The $3 \times 3$ minors, $4 \times 4$ minors, etc., impose increasingly
complex constraints on the interactions among groups of three, four, or more
lottery-parameter pairs. These are harder to interpret individually but
follow the same logic: no subset of lotteries should be ``too coupled'' to
each other's parameters relative to their own-parameter responses.

The final condition (order $p$) is simply $\det \widetilde{J} > 0$---the
full determinant of the unnormalized Jacobian must be positive.

\subsection{Diagonal Entries in Detail: When Are They Positive?}
\label{sec:diagonal_detail}

Since the diagonal entries are the most basic necessary condition for a
P-matrix, we examine what determines their sign.

\paragraph{Utility parameter diagonal: $A_{kk}$.}
From~\eqref{eq:A_entry}, specializing to the diagonal ($i = k$, meaning
lottery $k$ is assigned to utility parameter $\theta_k$):
\begin{equation}
    A_{kk} = \sum_{m} \pi_{mk} \,
    \frac{\partial u}{\partial \theta_k}(x_{mk}) -
    \frac{\partial u}{\partial \theta_k}(\CE_k)
\end{equation}

To determine the sign, consider the function
$h(x) = \partial u / \partial \theta_k(x;\, \btheta)$, viewed as a function
of $x$ alone (holding $\btheta$ fixed). Then:
\begin{equation}
    A_{kk} = \sum_m \pi_{mk}\, h(x_{mk}) - h(\CE_k)
    = \E_\pi[h(X)] - h(\CE_k)
\end{equation}
where $\E_\pi$ denotes the expectation under the decision weights.

Now, recall that the CE satisfies $u(\CE_k) = \E_\pi[u(X)]$. The CE is a
kind of ``weighted mean'' of the prizes, but in utility space rather than in
dollar space. If $h = \partial u / \partial \theta_k$ is convex in $x$, then
Jensen's inequality gives $\E_\pi[h(X)] > h(\CE_k)$, so $A_{kk} > 0$.
If $h$ is concave, then $A_{kk} < 0$.

The important point: $A_{kk} = 0$ requires $h$ to be \emph{affine} in $x$
on the prize range, which is a non-generic condition. For any utility
function where the marginal effect of $\theta_k$ on utility depends
nonlinearly on the outcome level $x$, the diagonal entry $A_{kk}$ is
generically nonzero. Its sign depends on the convexity/concavity of $h$, which
is a property of the specific utility function.

\paragraph{PWF parameter diagonal: $B_{kk}$.}
From~\eqref{eq:B_entry}, specializing to the diagonal ($j = k$, meaning
lottery $k$ is assigned to PWF parameter $\psi_k$):
\begin{equation}
    B_{kk} = \sum_{m} \frac{\partial \pi_{mk}}{\partial \psi_k} \,
    u(x_{mk})
\end{equation}

Since the weight perturbations sum to zero
($\sum_m \partial\pi_{mk}/\partial\psi_k = 0$), this is a \emph{mean-zero
reweighting} of the utility values. Think of it as asking: when $\psi_k$
increases, does it move decision weight toward the outcomes with
\emph{higher} utility, or toward those with \emph{lower} utility?

\begin{itemize}
    \item If $\psi_k$ shifts weight toward high-utility outcomes (the
    expensive prizes), then $B_{kk} > 0$: the CE increases because the
    lottery is now ``valued more generously.''

    \item If $\psi_k$ shifts weight toward low-utility outcomes (the cheap
    prizes), then $B_{kk} < 0$: the CE decreases.
\end{itemize}

The sign is determined by the interaction between the PWF's parametric
structure and the specific probability configuration of lottery $L_k$. Unlike
$A_{kk}$, there is no single universal sign---it depends on the model. The
researcher's job is to choose lottery designs (prize vectors and probability
structures) such that the diagonal entries come out with a consistent sign
that makes the P-matrix condition achievable.

\subsection{Diagonal Dominance: A Simpler Sufficient Condition}

Verifying all $2^p - 1$ principal minors individually can be tedious,
especially as $p$ grows. Fortunately, there is a simpler condition that
\emph{implies} the P-matrix property without checking each minor separately.

The idea is intuitive: if every lottery's response to its \emph{own}
parameter is much larger (in absolute value) than its combined responses to
all \emph{other} parameters, then the cross-effects cannot overwhelm the
own-effects at any order, and all principal minors are automatically positive.

\begin{proposition}[Diagonal dominance implies P-matrix]
\label{prop:diag_dom}
If $\widetilde{J}$ has positive diagonal entries and is strictly diagonally
dominant---meaning that for each row $k$, the diagonal entry exceeds the sum
of the absolute values of all off-diagonal entries:
\begin{equation}
    \widetilde{J}_{kk} > \sum_{\ell \neq k} |\widetilde{J}_{k\ell}|
    \qquad \text{for all } k = 1, \ldots, p
\end{equation}
---then $\widetilde{J}$ is a P-matrix.
\end{proposition}

\begin{proof}
A strictly diagonally dominant matrix with positive diagonal entries is a
(non-singular) M-matrix up to sign, and all its eigenvalues have positive real
parts. Every such matrix is a P-matrix (see, e.g., Horn and
Johnson~\cite{horn_johnson_2012}, Theorem 2.5.12, or Fiedler and
Pt\'ak~\cite{fiedler_ptak_1962}).
\end{proof}

In economic terms, diagonal dominance says: \emph{each lottery is most
informative about its own assigned parameter}. The CE of lottery $k$ moves
more when you change parameter $\xi_k$ than it does when you change all other
parameters combined. This is a natural design goal for the reference lotteries:
choose each lottery so that it is ``maximally sensitive'' to one particular
parameter.

Diagonal dominance is strictly stronger than the P-matrix condition---there
exist P-matrices that are not diagonally dominant. But it is much easier to
check: only $p$ row-sum conditions instead of $2^p - 1$ minor conditions.
For large $p$, this computational saving is decisive.

Diagonal dominance says: each lottery $k$ is \emph{most sensitive} to its own
assigned parameter $\xi_k$, in the sense that the own-effect dominates the sum
of all cross-effects. This is the formal version of the intuition that each
lottery should be ``most informative'' about one parameter.


%==============================================================================
\section{Prize Diversity and the P-Matrix}
\label{sec:prize_diversity}
%==============================================================================

This section identifies the economic mechanism---prize diversity---that
delivers the P-matrix property.

\subsection{The Shared-Prize Problem}

Suppose all $K$ lotteries use the \emph{same} prize vector
$(x_1, \ldots, x_M)$ but differ only in their probability structures. Then:
\begin{itemize}
    \item The utility values $u(x_m; \btheta)$ and their parameter derivatives
    $\partial u / \partial \theta_i(x_m; \btheta)$ are \emph{common} across
    all lotteries.

    \item The entries of the $A$-block become:
    \begin{equation}
        A_{ki} = \sum_{m=1}^M \pi_{mk} \,
        \frac{\partial u}{\partial \theta_i}(x_m)
        - \frac{\partial u}{\partial \theta_i}(\CE_k)
    \end{equation}
    The first sum, $\sum_m \pi_{mk} (\partial u / \partial \theta_i)(x_m)$,
    differs across lotteries $k$ only through the decision weights $\pi_{mk}$.
    The second term $(\partial u / \partial \theta_i)(\CE_k)$ differs across
    lotteries because different probability structures produce different CEs.
    However, both variations are driven by the \emph{same underlying function}
    $\partial u / \partial \theta_i$ evaluated at points within the common
    prize range $[x_1, x_M]$.
\end{itemize}

When there is only one utility parameter ($p_1 = 1$), the $A$-block is a
single column, so proportionality is not an issue. But the $A$-block's
contribution to the overall Jacobian structure is limited: all variation in
$A_{k1}$ across lotteries $k$ comes from how different probability structures
weight a single function $\partial u / \partial \theta_1$.

For the $B$-block under shared prizes: since $u(x_m; \btheta)$ is common
across lotteries, variation in $B_{kj}$ across $k$ comes entirely from
$\partial\pi_{mk}/\partial\psi_j$, which depends on the probability structure
of lottery $L_k$. If two lotteries have similar probability structures, their
rows in $B$ are nearly proportional, and $2 \times 2$ principal minors
involving those rows approach zero.

\begin{example}[Failure with shared prizes: the 3-parameter Prelec case]
\label{ex:shared_failure}
In the companion paper's application, $p_1 = 1$ (CRRA parameter $r$) and
$p_2 = 2$ (Prelec parameters $\phi, \eta$). Three lotteries on the same
prizes $(x_1, x_2, x_3)$ give a $3 \times 3$ Jacobian $\widetilde{J} =
[A \mid B]$ where $A$ is $3 \times 1$ and $B$ is $3 \times 2$.

The Prelec PWF $w(p) = \exp(-\eta(-\ln p)^\phi)$ has the property that
$\partial w / \partial \phi$ and $\partial w / \partial \eta$ are \emph{not}
proportional functions of $p$ (since
$(\partial w/\partial \phi) / (\partial w/\partial \eta) =
\eta \cdot (-\ln p)^\phi \cdot \ln(-\ln p)$ depends on $p$). This means the
$B$-block's two columns are generically linearly independent.

However, the numerical analysis in the companion code
(\texttt{injectivity\_analysis.py}) demonstrates that for three lotteries on
\emph{shared} prizes, distinct parameter triples $(r, \phi, \eta)$ and
$(r', \phi', \eta')$ can produce identical CE vectors. The $3 \times 3$
Jacobian is nonsingular at most points (local injectivity), but the mapping
folds---it is not globally injective. The P-matrix condition fails at some
points in the parameter space, and \Cref{thm:gale_nikaido} does not
certify injectivity.
\end{example}

\subsection{How Prize Diversity Restores the P-Matrix}

Now suppose the lotteries have \emph{distinct prize vectors}: lottery $L_k$
has prizes $(x_{1k}, \ldots, x_{M_k,k})$ that differ across $k$.

\paragraph{Effect on the $A$-block.}
With distinct prizes, the function $\partial u / \partial \theta_i$ is
evaluated at \emph{different points} for different lotteries. The entries of
column $i$ of $A$ become:
\begin{equation}
    A_{ki} = \sum_{m} \pi_{mk} \,
    \frac{\partial u}{\partial \theta_i}(x_{mk})
    - \frac{\partial u}{\partial \theta_i}(\CE_k)
\end{equation}
where the arguments $x_{mk}$ now vary with $k$. This gives each row of $A$
its own ``fingerprint'' determined by the specific prizes of lottery $L_k$,
breaking any proportionality between rows.

\paragraph{Effect on the $B$-block.}
With distinct prizes:
\begin{equation}
    B_{kj} = \sum_{m}
    \frac{\partial \pi_{mk}}{\partial \psi_j} \, u(x_{mk})
\end{equation}
The utility values $u(x_{mk})$ now differ across lotteries, so even if
$\partial\pi_{mk}/\partial\psi_j$ were similar across lotteries (which it
generically is not, since probability structures also differ), the different
utility values prevent proportionality of rows.

\paragraph{Effect on cross-block interactions.}
The $2 \times 2$ and higher-order principal minors involve entries from both
$A$ and $B$. Prize diversity ensures that the mixed minors
$A_{ki} B_{\ell j} - A_{\ell i} B_{kj}$ are generically nonzero, because the
terms depend on prizes specific to lotteries $k$ and $\ell$.

\begin{keyresult}[Prize diversity mechanism]
Prize diversity delivers the P-matrix property through two channels:
\begin{enumerate}
    \item \textbf{Row differentiation:} Distinct prizes give each lottery its
    own utility gradient profile, breaking proportionality between rows of
    $\widetilde{J}$ (whether within $A$, within $B$, or across blocks).

    \item \textbf{Approximate diagonal dominance:} With sufficiently diverse
    prizes, each lottery can be designed so that its CE is ``most sensitive'' to
    one designated parameter, creating approximate diagonal dominance of
    $\widetilde{J}$.
\end{enumerate}
\end{keyresult}


%==============================================================================
\section{Formal Sufficient Conditions}
\label{sec:main_theorem}
%==============================================================================

We now state the main result.

\begin{theorem}[Global injectivity of the CE mapping]
\label{thm:main}
Let $u(x; \btheta)$ and $w(p; \bpsi)$ satisfy the primitives and
monotone-influence assumptions of Section~\ref{sec:model}. Let
$\Xi = \Theta \times \Psi$ be a rectangular domain (product of open
intervals). Let $L_1, \ldots, L_K$ be $K = p_1 + p_2$ reference lotteries,
and define the CE mapping $g: \Xi \to \R^K$ by~\eqref{eq:g_map}. If:
\begin{enumerate}
    \item[(i)] \textbf{Rectangular domain:} $\Xi = \prod_{i=1}^p (a_i, b_i)$
    with $-\infty \leq a_i < b_i \leq +\infty$;

    \item[(ii)] \textbf{Monotone influence:}
    $\partial u / \partial \theta_i \neq 0$ and
    $\partial w / \partial \psi_j \neq 0$ on $\Xi$ for all $i, j$;

    \item[(iii)] \textbf{P-matrix Jacobian:}
    $\widetilde{J}(\bxi)$ is a P-matrix for all $\bxi \in \Xi$;
\end{enumerate}
then $g$ is globally injective on $\Xi$.
\end{theorem}

\begin{proof}
By \Cref{lem:D_reduction}, condition (iii) is equivalent to $J_g(\bxi)$
being a P-matrix for all $\bxi \in \Xi$. Since $\Xi$ is rectangular by
condition (i), the Gale--Nikaido theorem (\Cref{thm:gale_nikaido})
applies directly, giving global injectivity of $g$ on $\Xi$.
\end{proof}

\begin{remark}[Roles of the conditions]
Conditions (i) and (ii) are \emph{model properties}: the domain is rectangular
by construction (each parameter has its own range), and monotone influence is
a minimal requirement for each parameter to be identifiable. The operative
condition is (iii)---the P-matrix property---which depends on the
\emph{interaction} between the model, the parameter values, and the lottery
designs.

Prize diversity (discussed in Section~\ref{sec:prize_diversity}) is an
instrument for achieving condition (iii): it is neither necessary nor
sufficient in the abstract, but it is the economic mechanism that
\emph{generically} delivers the P-matrix property. With sufficient prize
diversity and well-separated probability structures, the Jacobian
$\widetilde{J}$ is approximately diagonally dominant, hence a P-matrix.
\end{remark}

\begin{remark}[The role of lottery design]
The researcher controls the choice of reference lotteries $L_1, \ldots, L_K$.
This is a degree of freedom not present in most applications of the
Gale--Nikaido theorem (where the mapping is given). The implication is that
\emph{lottery design} is a tool for ensuring global injectivity: by choosing
lotteries with appropriate prize vectors and probability structures, the
researcher can engineer the Jacobian to satisfy the P-matrix condition.
\end{remark}

\subsection{The Overdetermined Case ($K > p$)}
\label{sec:overdetermined}

The main theorem (\Cref{thm:main}) treats the square case $K = p$. In
practice, experimenters often have $K > p$ reference lotteries---more
equations than parameters. This overdetermined case deserves separate
treatment because the Jacobian is no longer square and the P-matrix concept
does not directly apply.

\paragraph{The non-square Jacobian.}
When $K > p$, the Jacobian $J_g(\bxi) \in \R^{K \times p}$ is a tall matrix
(more rows than columns). P-matrices are defined only for square matrices
(principal minors require selecting the \emph{same} row and column indices),
so one cannot check whether $J_g$ is a P-matrix. The Gale--Nikaido theorem
does not apply to $J_g$ directly.

\paragraph{The Gram matrix $G = J_g^\top J_g$.}
A natural candidate for a square matrix is the Gram matrix
$G(\bxi) = J_g(\bxi)^\top J_g(\bxi) \in \R^{p \times p}$. This matrix
arises in least-squares analysis and determines local identification: $G$
is positive definite if and only if $J_g$ has full column rank $p$, which
is equivalent to local injectivity at $\bxi$.

However, the Gram matrix does \textbf{not} plug into the Gale--Nikaido
theorem for global injectivity. The theorem requires the Jacobian of the
mapping itself to be a P-matrix---not a derived quantity. The Gram matrix
$G = J_g^\top J_g$ is positive semidefinite by construction, and its
principal minors are nonneg\-ative, but:
\begin{enumerate}
    \item $G$ is positive semidefinite, not necessarily a P-matrix (P-matrix
    requires \emph{strict} positivity of all principal minors, plus a
    structural condition beyond mere positive definiteness---though in fact,
    every symmetric positive definite matrix \emph{is} a P-matrix).

    \item More fundamentally, $G$ is not the Jacobian of any mapping
    $\Xi \to \R^p$. The Gale--Nikaido theorem applies to
    $J_{\mathbf{f}}(\bxi)$ where $\mathbf{f}: \R^p \to \R^p$ is the mapping
    whose injectivity is in question. One cannot substitute an arbitrary
    $p \times p$ matrix derived from the problem and invoke the theorem.
\end{enumerate}

\paragraph{Strategy 1: Extract a $p \times p$ sub-Jacobian.}
The most direct approach is to \emph{select} $p$ of the $K$ lotteries to form
a square subsystem. If lotteries $L_{k_1}, \ldots, L_{k_p}$ are selected,
define the restricted mapping $g_S: \Xi \to \R^p$ by
$g_S(\bxi) = (\CE_{k_1}(\bxi), \ldots, \CE_{k_p}(\bxi))$. Its Jacobian
$J_{g_S}(\bxi)$ is the $p \times p$ submatrix of $J_g$ formed by selecting
rows $k_1, \ldots, k_p$.

If $J_{g_S}(\bxi)$ is a P-matrix for all $\bxi \in \Xi$, then $g_S$ is
globally injective on $\Xi$ by \Cref{thm:main}. Since $g_S$ injective
implies $g$ injective (if the restricted mapping is one-to-one, the full
mapping with additional components is \emph{a fortiori} one-to-one), this
gives:

\begin{corollary}[Overdetermined case via subsystem selection]
\label{cor:overdetermined_sub}
If there exists a subset $S \subseteq \{1, \ldots, K\}$ with $|S| = p$ such
that the $p \times p$ sub-Jacobian $J_{g_S}(\bxi)$ is a P-matrix for all
$\bxi \in \Xi$, then the full CE mapping $g: \Xi \to \R^K$ is globally
injective on $\Xi$.
\end{corollary}

The researcher's task is then to find a ``good'' subset $S$---one for which
the P-matrix condition holds. Prize diversity helps here: with $K > p$
lotteries, there are $\binom{K}{p}$ possible subsets, and it suffices for
\emph{any one} of them to yield a P-matrix Jacobian.

\paragraph{Strategy 2: Generic injectivity from overdetermination.}
The overdetermined case also admits a simpler, more general argument.
When $K > p$, the system of equations $g(\bxi) = \mathbf{c}$ (for a target CE
vector $\mathbf{c} \in \R^K$) is overconstrained: $K$ equations in $p$
unknowns. Generically, such a system has at most one solution---the extra
equations eliminate the ``compensation'' degrees of freedom that allow
non-injectivity in the square case. More precisely:

\begin{proposition}[Generic injectivity from overdetermination]
\label{prop:overdetermined_generic}
Suppose $J_g(\bxi)$ has full column rank $p$ for all $\bxi \in \Xi$ (i.e.,
the Gram matrix $G(\bxi) = J_g^\top J_g$ is positive definite everywhere).
If $K > p$, then for a generic choice of $K$ lotteries, $g$ is globally
injective on $\Xi$.
\end{proposition}

The word ``generic'' here means: the set of lottery designs for which $g$
fails to be injective has measure zero in the space of possible lottery
designs. This is a transversality argument---adding lotteries beyond $K = p$
imposes additional constraints that generically prevent the coincidences
needed for non-injectivity.

\begin{remark}[When the Gram matrix is informative]
\label{rem:gram_informative}
Although $G = J_g^\top J_g$ does not feed into Gale--Nikaido, it plays two
useful roles in the overdetermined case:
\begin{enumerate}
    \item \textbf{Local identification:} $G(\bxi)$ positive definite
    (equivalently, $J_g$ has full column rank) is necessary and sufficient
    for local injectivity at $\bxi$. If $G$ is singular at some $\bxi_0$,
    there exists a direction $v \in \R^p$ with $J_g v = 0$, meaning the CE
    vector does not change (to first order) when parameters move in direction
    $v$---a failure of local identification.

    \item \textbf{Condition number:} The condition number
    $\kappa(G) = \lambda_{\max}(G) / \lambda_{\min}(G)$ measures how
    ``close to singular'' the Jacobian is. A large $\kappa$ indicates that
    some parameter combinations are poorly identified by the CE
    mapping---even if technically injective, small CE differences correspond
    to large parameter differences, making the inversion numerically
    ill-conditioned.
\end{enumerate}
\end{remark}

\paragraph{Comparison of approaches.}
\Cref{tab:square_vs_over} summarizes the tools available for each case.

\begin{table}[ht]
\centering
\begin{tabular}{lcc}
    & $K = p$ (square) & $K > p$ (overdetermined) \\
    \hline
    Jacobian shape & $p \times p$ & $K \times p$ \\
    P-matrix on $J_g$? & Yes (Gale--Nikaido) & Not defined \\
    Gram matrix $J_g^\top J_g$ PD? & $\Leftrightarrow \det J_g \neq 0$
        & Full rank (local ID) \\
    Global injectivity tool & P-matrix $\Rightarrow$ injective
        & Subsystem P-matrix, or generic argument \\
    Prize diversity role & Delivers P-matrix & Delivers full rank +
        overconstrained
\end{tabular}
\caption{Tools for global injectivity in the square vs.\ overdetermined case.}
\label{tab:square_vs_over}
\end{table}

\paragraph{Practical recommendation.}
When the researcher has $K > p$ reference lotteries (as is typical in
experimental economics, where task counts far exceed parameter counts), the
most robust strategy is:
\begin{enumerate}
    \item Verify $J_g$ has full column rank everywhere on the prior support
    (i.e., $G = J_g^\top J_g$ is positive definite)---this confirms local
    injectivity.
    \item Select a subset $S$ of $p$ lotteries with maximal prize diversity
    and verify the P-matrix condition on the $p \times p$ sub-Jacobian
    $J_{g_S}$---this gives a rigorous Gale--Nikaido certificate for global
    injectivity.
    \item If the P-matrix check is infeasible, the combination of full column
    rank and overdetermination provides strong (generically sufficient)
    evidence for global injectivity.
\end{enumerate}


%==============================================================================
\section{Practical Verification}
\label{sec:verification}
%==============================================================================

Before discussing verification methods, it is worth emphasizing a fundamental
point:

\begin{keyresult}[Verification is a one-time cost]
The P-matrix condition depends on the \emph{model specification} (utility
function, PWF) and the \emph{reference lottery designs} (prize vectors,
probability structures). It does \textbf{not} depend on the data, the prior
hyperparameters, the number of subjects, or any other aspect of a specific
empirical application.

Once the P-matrix property has been established for a given model and lottery
design, it holds for \emph{every} application using that model and those
lotteries---regardless of the dataset, the sample size, or the prior
specification. The verification cost is paid once per model-lottery
combination, not once per estimation run or once per MCMC iteration.
\end{keyresult}

This means that even a computationally expensive verification (symbolic
analysis, dense grid search, interval arithmetic) is entirely practical: the
cost is amortized over all future applications of the model. For commonly
used models like CRRA + Prelec, a single verification effort produces a
permanent result that subsequent researchers can cite without re-deriving.

\subsection{Enumerating Principal Minors}

For a $p \times p$ matrix, there are $2^p - 1$ principal minors.
\Cref{tab:minor_count} shows the count for small $p$.

\begin{table}[ht]
\centering
\begin{tabular}{ccc}
    $p$ & $2^p - 1$ & Breakdown by size \\
    \hline
    2 & 3 & $2 + 1$ \\
    3 & 7 & $3 + 3 + 1$ \\
    4 & 15 & $4 + 6 + 4 + 1$ \\
    5 & 31 & $5 + 10 + 10 + 5 + 1$ \\
    10 & 1023 & $10 + 45 + 120 + \cdots + 1$
\end{tabular}
\caption{Number of principal minors for the P-matrix condition.}
\label{tab:minor_count}
\end{table}

For the companion paper's 3-parameter case ($p = 3$), there are only
7 principal minors:
\begin{itemize}
    \item 3 diagonal entries (order 1)
    \item 3 off-diagonal $2 \times 2$ minors (order 2)
    \item 1 full determinant (order 3)
\end{itemize}

\subsection{Symbolic/Analytic Verification}
\label{sec:symbolic}

For any specific model (specific functional forms for $u$ and $w$), the
strongest approach is to derive the principal minors \emph{symbolically} and
prove analytically that they are positive on the parameter domain. A computer
algebra system (CAS) such as SymPy, Mathematica, or Maple can:

\begin{enumerate}
    \item Compute the entries of $\widetilde{J}(\bxi)$ symbolically from the
    closed-form expressions for $u(x;\btheta)$ and $w(p;\bpsi)$.

    \item Extract all $2^p - 1$ principal minors as symbolic expressions in
    $\bxi$.

    \item Simplify each minor and determine its sign on the parameter domain.
    For many utility/PWF combinations, the minors factor into products of
    terms whose signs can be established individually (e.g., sums of positive
    quantities, or terms that are manifestly positive by the monotonicity
    assumptions on $u$ and $w$).

    \item If a minor does not have a transparent sign, the CAS can search for
    zeros: solve $M_I(\bxi) = 0$ for each minor $M_I$ and check whether any
    solutions lie within the parameter domain $\Xi$. If none do, the minor is
    sign-definite on $\Xi$. Tools like Gr\"obner bases and cylindrical
    algebraic decomposition can handle this for polynomial or rational
    expressions; for transcendental functions (like the Prelec PWF), numerical
    root-finding over the bounded prior support can provide strong evidence.
\end{enumerate}

This approach yields a \emph{proof} (or, for transcendental cases, a
near-proof backed by exhaustive numerical search) rather than merely
numerical evidence. For small $p$---and $p = 3$ to $5$ covers most
individual-level models in behavioral economics---the symbolic expressions
are manageable.

\begin{remark}[Existing symbolic infrastructure]
\label{rem:sympy}
The companion code (\texttt{injectivity\_analysis.py}, Analysis 1) already
computes the Jacobian $\widetilde{J}$ symbolically using SymPy for the
CRRA + Prelec case, including numerical evaluation of $\det \widetilde{J}$
at test points. Extending this to compute and analyze all 7 principal minors
(rather than just the full determinant) is straightforward: extract the three
$1 \times 1$ and three $2 \times 2$ principal submatrices from the symbolic
Jacobian and compute their determinants.
\end{remark}

\subsection{Numerical Grid Evaluation}

When symbolic analysis is intractable (e.g., for complex utility functions
that resist closed-form simplification), the most direct fallback is to
evaluate all $2^p - 1$ principal minors of $\widetilde{J}(\bxi)$ on a grid
covering the prior support $\Xi_0 \subset \Xi$:

\begin{enumerate}
    \item Define a grid $\{\bxi^{(1)}, \ldots, \bxi^{(N)}\}$ over $\Xi_0$
    with spacing fine enough to detect sign changes.

    \item At each grid point $\bxi^{(n)}$:
    \begin{enumerate}
        \item Compute $\widetilde{J}(\bxi^{(n)})$ (via automatic
        differentiation or closed-form expressions).
        \item Evaluate all $2^p - 1$ principal minors.
        \item Record the minimum minor value.
    \end{enumerate}

    \item If all minors are positive at all grid points, this provides strong
    numerical evidence for the P-matrix property on $\Xi_0$.
\end{enumerate}

\paragraph{Computational cost.} For a grid with $N$ points and $p$ parameters,
the cost is $O(N \cdot 2^p \cdot p^3)$ (evaluating $2^p - 1$ determinants of
submatrices up to size $p \times p$). For $p = 3$ and $N = 10{,}000$, this
is $7 \times 10{,}000 = 70{,}000$ determinant evaluations, each of at most
$3 \times 3$ matrices---a computation that takes seconds on a modern machine.
Even for a dense grid with $N = 10^6$ points, the total is $7 \times 10^6$
small determinants, which remains trivial.

Since this is a one-time computation (see the key result box above), even an
expensive grid search is entirely acceptable. One can afford to be thorough:
use a very dense grid, cover a generous parameter range well beyond the prior
support, and store the results for future reference.

\subsection{Diagonal Dominance Check}

When the full P-matrix verification is unwieldy (very large $p$, or when one
wants a quick sanity check before investing in a full analysis), diagonal
dominance (\Cref{prop:diag_dom}) provides a sufficient condition requiring
only $p$ row-sum checks:
\begin{equation}
    \widetilde{J}_{kk} > \sum_{\ell \neq k}
    |\widetilde{J}_{k\ell}|, \quad k = 1, \ldots, p
\end{equation}
This is $O(N \cdot p^2)$ per grid point---polynomial rather than exponential
in $p$. However, diagonal dominance is a \emph{stronger} condition than
P-matrix, so it may fail even when the P-matrix property holds. It is best
used as a first-pass screening: if it succeeds, the verification is complete;
if it fails, one proceeds to the full minor analysis.

\subsection{Interval Arithmetic for Rigorous Certification}

Grid evaluation provides strong evidence but not a proof: in principle, the
minors could change sign between grid points. For rigorous certification
(e.g., for a result stated as a theorem in a published paper), \emph{interval
arithmetic} eliminates this gap:
\begin{enumerate}
    \item Partition $\Xi_0$ into boxes $B_1, \ldots, B_N$.
    \item On each box $B_n$, compute an interval enclosure of each
    principal minor using interval-arithmetic evaluation of $\widetilde{J}$.
    \item If every interval enclosure has a strictly positive lower bound,
    the P-matrix property is rigorously certified on $\Xi_0$.
\end{enumerate}

This approach requires an interval-arithmetic implementation of the CE
mapping and its Jacobian. Tools such as \texttt{mpmath} (Python) or
\texttt{IntervalArithmetic.jl} (Julia) provide the necessary infrastructure.

Again, since verification is one-time, the computational cost of interval
arithmetic---which is typically $10$--$100\times$ that of ordinary
floating-point evaluation---is entirely acceptable.

\subsection{Recommended Verification Strategy}

For a specific model and lottery design, we recommend the following sequence:

\begin{enumerate}
    \item \textbf{Symbolic analysis (preferred).} Use a CAS to derive the
    principal minors symbolically. If their positivity can be established
    analytically (or if the CAS can certify no zeros in the domain), this
    provides a complete proof. For the typical $p = 3$ to $5$ case in
    behavioral economics, this is the most productive approach.

    \item \textbf{Dense grid evaluation (fallback).} If symbolic analysis
    stalls, evaluate all minors on a dense grid over the prior support.
    Record the minimum minor value as a function of the parameter vector.
    This is cheap (one-time, seconds to minutes) and provides strong
    evidence.

    \item \textbf{Interval arithmetic (certification).} If a formal proof is
    needed and symbolic methods did not yield one, use interval arithmetic to
    rigorously bound the minors from below on the prior support.
\end{enumerate}

At each stage, the result is permanent: once established for a given
model-lottery pair, it applies to all future applications.


%==============================================================================
\section{Application: CRRA + Prelec (the Paper's Case)}
\label{sec:application}
%==============================================================================

We specialize the general framework to the companion paper's model:
$p_1 = 1$ (CRRA parameter $r$), $p_2 = 2$ (Prelec parameters $\phi, \eta$),
$K = 3$ reference lotteries.

\subsection{Model Specification}

\begin{itemize}
    \item \textbf{Utility:} $u(x; r) = x^{1-r} / (1-r)$ for $r \neq 1$,
    $\ln x$ for $r = 1$.

    \item \textbf{PWF:} $w(p; \phi, \eta) = \exp(-\eta(-\ln p)^\phi)$
    (Prelec 2-parameter).

    \item \textbf{Parameter domain:}
    $\Xi = (r_{\min}, r_{\max}) \times (\phi_{\min}, \phi_{\max})
    \times (\eta_{\min}, \eta_{\max})$, a rectangular domain.
\end{itemize}

The monotone-influence conditions are satisfied:
\begin{itemize}
    \item $\partial u / \partial r = -x^{1-r} \ln x / (1-r)
    + x^{1-r} / (1-r)^2 \neq 0$ for $x > 0$ (can be verified to be
    nonvanishing on the relevant domain, or one works directly with
    the $B_k = (1-r) \cdot u(\CE_k; r)$ formulation where
    $\partial B_k / \partial r$ is explicit).

    \item $\partial w / \partial \phi =
    -w(p) \cdot \eta \cdot (-\ln p)^\phi \cdot \ln(-\ln p) \neq 0$ for
    $p \in (0, 1) \setminus \{e^{-1}\}$.

    \item $\partial w / \partial \eta =
    -w(p) \cdot (-\ln p)^\phi \neq 0$ for $p \in (0, 1)$.
\end{itemize}

\subsection{Lottery Design}
\label{sec:lottery_design}

We use $K = 3$ reference lotteries with distinct prize vectors.
Each lottery is a three-outcome gamble with prizes $x_1 < x_2 < x_3$
and probabilities $(p_1, p_2, p_3)$.
Define cumulative probabilities from above:
$q \coloneqq p_2 + p_3$ (probability of receiving $x_2$ or better)
and $s_{\mathrm{cum}} \coloneqq p_3$ (probability of receiving $x_3$).
The three decision weights under RDU are:
\begin{align}
    \pi_1 &= 1 - w(q), \qquad
    \pi_2 = w(q) - w(s_{\mathrm{cum}}), \qquad
    \pi_3 = w(s_{\mathrm{cum}})
\end{align}

The specific design assigns one lottery to each parameter
$(r, \phi, \eta)$:
\begin{center}
\begin{tabular}{lccccl}
    Lottery & Prizes & $q$ & $s_{\mathrm{cum}}$ & Assigned to & Sign of $\partial\CE/\partial\xi$ \\
    \hline
    $L_1$ & $(10, 30, 50)$ & $0.50$ & $0.50$ & $r$ & $-$ \\
    $L_2$ & $(15, 35, 50)$ & $0.60$ & $0.90$ & $\phi$ & $+$ \\
    $L_3$ & $(10, 20, 50)$ & $0.15$ & $0.25$ & $\eta$ & $-$
\end{tabular}
\end{center}

\noindent
Note that $L_1$ has $q = s_{\mathrm{cum}} = 0.50$, so $\pi_2 = w(0.50) - w(0.50) = 0$:
lottery $L_1$ is effectively a \emph{binary} gamble on $\{10, 50\}$ with
decision weights $\pi_1 = 1 - w(0.50)$ and $\pi_3 = w(0.50)$.
This simplification will be useful below.

The ``Sign'' column records the generic sign of $\partial\CE_k/\partial\xi_k$
(the diagonal entry of the Jacobian for the assigned parameter).
Since the P-matrix condition requires positive diagonal entries, we
apply a \emph{sign reparameterization}: define $S = \diag(-1, +1, -1)$
and consider the modified mapping $\tilde{g} = S \circ g$. Since $S$ is
a fixed sign matrix, $g$ is injective if and only if $\tilde{g}$ is
injective. The Jacobian of $\tilde{g}$ is $S \cdot J_g$, which negates
the $r$ and $\eta$ columns of $J_g$. After this reparameterization, all
diagonal entries of $\widetilde{J}$ are positive, as required.

\subsection{The Power-Mean Representation}
\label{sec:power_mean}

For CRRA utility $u(x;r) = x^s/s$ where $s \coloneqq 1 - r$
(with the convention $u = \ln x$ at $s = 0$), the CE has a
closed-form \emph{power-mean} representation.
Define the \textbf{power-mean body}:
\begin{equation}
    \label{eq:B_def}
    B_k \;\coloneqq\; \sum_{m=1}^{M_k} \pi_{mk}(\phi,\eta)\; x_{mk}^{\,s}
\end{equation}
Then the RDU value is $\RDU_k = B_k / s$ and the certainty equivalent is:
\begin{equation}
    \label{eq:CE_power}
    \CE_k \;=\; B_k^{\,1/s}
\end{equation}
This is simply the \emph{weighted power mean} of the prizes under
decision weights $\pi_{mk}$.

Because $B_k$ is a positive linear combination of $x_{mk}^s > 0$
(since decision weights are positive), $B_k > 0$ for all parameter values.
The CE is well-defined and positive.

\subsection{The $3 \times 3$ Jacobian in Explicit Form}
\label{sec:explicit_jacobian}

To apply the Gale--Nikaido theorem, we need the Jacobian of the
CE mapping. With the power-mean representation, we can express the
derivatives of $\CE_k = B_k^{1/s}$ in terms of $B_k$ and its partial
derivatives.

\paragraph{Derivative with respect to $r$.}
Since $s = 1-r$, we have $\partial/\partial r = -\partial/\partial s$.
By logarithmic differentiation of $\CE_k = B_k^{1/s}$:
\begin{equation}
    \label{eq:dCE_dr}
    \frac{\partial \CE_k}{\partial r}
    \;=\; \frac{\CE_k}{s}
    \left[
        \frac{\ln B_k}{s}
        - \frac{1}{B_k}\sum_{m} \pi_{mk}\, x_{mk}^s \ln x_{mk}
    \right]
\end{equation}
The bracketed expression can be rewritten as:
\begin{equation}
    \frac{\partial \CE_k}{\partial r}
    \;=\; \frac{\CE_k}{s}
    \Big[\ln(\CE_k) - \E_\pi[\ln X_k]\Big]
\end{equation}
where $\E_\pi[\ln X_k] \coloneqq (1/B_k)\sum_m \pi_{mk}\, x_{mk}^s \ln x_{mk}$
is a power-weighted average of $\ln x_{mk}$.
By Jensen's inequality applied to the concave function $\ln$, we have
$\ln(\CE_k) > \E_\pi[\ln X_k]$ (since the prizes are not all equal),
so the bracket is positive. Combined with $\CE_k > 0$, the sign of
$\partial\CE_k/\partial r$ is determined by the sign of $1/s = 1/(1-r)$:
negative for $r < 1$, positive for $r > 1$.
On any domain excluding $r = 1$ (a single point of Lebesgue
measure zero where CRRA becomes $\ln x$), this derivative is
sign-definite.

\paragraph{Derivatives with respect to $\phi$ and $\eta$.}
Since $\phi$ and $\eta$ enter only through the decision weights
$\pi_{mk}$, the chain rule gives:
\begin{equation}
    \label{eq:dCE_dpsi}
    \frac{\partial \CE_k}{\partial \psi_j}
    \;=\; \frac{\CE_k}{s \cdot B_k}
    \;\frac{\partial B_k}{\partial \psi_j},
    \qquad \psi_j \in \{\phi, \eta\}
\end{equation}
where:
\begin{equation}
    \label{eq:dB_dpsi}
    \frac{\partial B_k}{\partial \psi_j}
    \;=\; \sum_{m=1}^{M_k}
    \frac{\partial \pi_{mk}}{\partial \psi_j}\; x_{mk}^{\,s}
\end{equation}
The prefactor $\CE_k / (s \cdot B_k)$ is common to both PWF derivatives
and is sign-definite (positive for $s > 0$, negative for $s < 0$).
Therefore the sign of $\partial\CE_k/\partial\psi_j$ is determined
entirely by $\partial B_k / \partial \psi_j$.

\paragraph{Factorization.}
From the general theory (Section~\ref{sec:jacobian}), the CE Jacobian
factors as $J_g = D^{-1}\widetilde{J}$ with $D$ a positive diagonal
matrix. After the sign reparameterization $S = \diag(-1, +1, -1)$
from Section~\ref{sec:lottery_design}, we need $S \cdot J_g$ to be a
P-matrix. Absorbing the positive diagonal $D^{-1}$ (which does not
affect the P-matrix property by \Cref{lem:D_reduction}), it suffices
to check that the \textbf{sign-adjusted unnormalized Jacobian}
\begin{equation}
    \label{eq:Jtilde_SA}
    \widetilde{J}^{\,\mathrm{SA}}
    \;\coloneqq\;
    \begin{pmatrix}
        -\widetilde{J}_{1,r}  & \widetilde{J}_{1,\phi}  & -\widetilde{J}_{1,\eta} \\[3pt]
        -\widetilde{J}_{2,r}  & \widetilde{J}_{2,\phi}  & -\widetilde{J}_{2,\eta} \\[3pt]
        -\widetilde{J}_{3,r}  & \widetilde{J}_{3,\phi}  & -\widetilde{J}_{3,\eta}
    \end{pmatrix}
\end{equation}
is a P-matrix for all $(r, \phi, \eta) \in \Xi$.

\subsection{Explicit Entries of $\widetilde{J}^{\,\mathrm{SA}}$ for the Lottery Design}

We now write out each entry of $\widetilde{J}^{\,\mathrm{SA}}$ for the
specific lottery design from Section~\ref{sec:lottery_design}.
Throughout, $s = 1 - r$, and we write
$t_p \coloneqq -\ln p$ and $w_p \coloneqq w(p;\phi,\eta) =
\exp(-\eta\, t_p^{\,\phi})$.

\paragraph{Prelec weight derivatives.}
The decision weights $\pi_{mk}$ are differences of Prelec
evaluations at cumulative probabilities. Their derivatives with
respect to $\phi$ and $\eta$ follow from:
\begin{align}
    \label{eq:dw_dphi}
    \frac{\partial w}{\partial \phi}(p)
    &= -w(p)\;\eta\; t_p^{\,\phi}\;\ln t_p \\[4pt]
    \label{eq:dw_deta}
    \frac{\partial w}{\partial \eta}(p)
    &= -w(p)\; t_p^{\,\phi}
\end{align}
Note the negative sign in~\eqref{eq:dw_dphi}: it arises from
$\frac{d}{d\phi}\exp(-\eta\, t_p^{\,\phi})
= \exp(-\eta\, t_p^{\,\phi}) \cdot (-\eta)\cdot t_p^{\,\phi}\ln t_p$.

\paragraph{Notation for the entries.}
For each lottery $L_k$, define the power-mean body $B_k$ as
in~\eqref{eq:B_def}. The sign-adjusted unnormalized Jacobian entries
are (up to positive row-dependent prefactors that do not affect the
P-matrix property):
\begin{itemize}
    \item \textbf{$r$-column} (negated, from $-\partial\CE_k/\partial r$):
    proportional to $\sum_m \pi_{mk}\,x_{mk}^s\,\ln x_{mk} - B_k\ln(\CE_k)$

    \item \textbf{$\phi$-column} (from $\partial\CE_k/\partial\phi$):
    proportional to $\partial B_k / \partial\phi$

    \item \textbf{$\eta$-column} (negated, from $-\partial\CE_k/\partial\eta$):
    proportional to $-\partial B_k / \partial\eta$
\end{itemize}

\paragraph{Row 1: Lottery $L_1$ (binary gamble on $\{10, 50\}$).}

Since $q = s_{\mathrm{cum}} = 0.50$, the middle outcome has zero
weight ($\pi_2 = 0$), giving:
\begin{equation}
    B_1 = (1 - w_{0.50})\cdot 10^s + w_{0.50}\cdot 50^s
\end{equation}
where $w_{0.50} \coloneqq w(0.50;\phi,\eta)
= \exp(-\eta\,(\ln 2)^\phi)$.
The derivative contributions simplify to two-term sums:
\begin{align}
    \frac{\partial B_1}{\partial\phi}
    &= \frac{\partial w_{0.50}}{\partial\phi}\;(50^s - 10^s)
    = -w_{0.50}\;\eta\;(\ln 2)^\phi\;\ln(\ln 2)\;(50^s - 10^s) \\[4pt]
    \frac{\partial B_1}{\partial\eta}
    &= \frac{\partial w_{0.50}}{\partial\eta}\;(50^s - 10^s)
    = -w_{0.50}\;(\ln 2)^\phi\;(50^s - 10^s)
\end{align}
Since $50^s > 10^s$ for $s > 0$ (and $50^s < 10^s$ for $s < 0$), the
factor $(50^s - 10^s)$ has the sign of $s$.

\paragraph{Row 2: Lottery $L_2$ (prizes $15, 35, 50$; $q = 0.60$,
$s_{\mathrm{cum}} = 0.90$).}

The cumulative probabilities are $q = 0.60$ and $s_{\mathrm{cum}} = 0.90$,
giving decision weights:
\begin{align}
    \pi_{1,2} &= 1 - w(0.60), \qquad
    \pi_{2,2} = w(0.60) - w(0.90), \qquad
    \pi_{3,2} = w(0.90)
\end{align}
The power-mean body is:
\begin{equation}
    B_2 = (1 - w_{0.60})\cdot 15^s
    + (w_{0.60} - w_{0.90})\cdot 35^s
    + w_{0.90}\cdot 50^s
\end{equation}
Its derivatives involve two PWF evaluation points:
\begin{align}
    \frac{\partial B_2}{\partial\phi}
    &= -\frac{\partial w_{0.60}}{\partial\phi}\;(15^s - 35^s)
    \;-\; \frac{\partial w_{0.90}}{\partial\phi}\;(35^s - 50^s) \\[4pt]
    \frac{\partial B_2}{\partial\eta}
    &= -\frac{\partial w_{0.60}}{\partial\eta}\;(15^s - 35^s)
    \;-\; \frac{\partial w_{0.90}}{\partial\eta}\;(35^s - 50^s)
\end{align}
where $\partial w_p/\partial\phi$ and $\partial w_p/\partial\eta$
are given by~\eqref{eq:dw_dphi}--\eqref{eq:dw_deta}.

\paragraph{Row 3: Lottery $L_3$ (prizes $10, 20, 50$; $q = 0.15$,
$s_{\mathrm{cum}} = 0.25$).}

Decision weights:
\begin{align}
    \pi_{1,3} &= 1 - w(0.15), \qquad
    \pi_{2,3} = w(0.15) - w(0.25), \qquad
    \pi_{3,3} = w(0.25)
\end{align}
Power-mean body:
\begin{equation}
    B_3 = (1 - w_{0.15})\cdot 10^s
    + (w_{0.15} - w_{0.25})\cdot 20^s
    + w_{0.25}\cdot 50^s
\end{equation}
Derivatives:
\begin{align}
    \frac{\partial B_3}{\partial\phi}
    &= -\frac{\partial w_{0.15}}{\partial\phi}\;(10^s - 20^s)
    \;-\; \frac{\partial w_{0.25}}{\partial\phi}\;(20^s - 50^s) \\[4pt]
    \frac{\partial B_3}{\partial\eta}
    &= -\frac{\partial w_{0.15}}{\partial\eta}\;(10^s - 20^s)
    \;-\; \frac{\partial w_{0.25}}{\partial\eta}\;(20^s - 50^s)
\end{align}

\paragraph{The $r$-column.}
For each lottery, the $r$-column entry (after sign negation) is
proportional to:
\begin{equation}
    \label{eq:r_col}
    R_k \;\coloneqq\;
    \sum_m \pi_{mk}\,x_{mk}^s\,\ln x_{mk}
    \;-\; B_k\,\ln(B_k^{1/s})
    \;=\;
    \sum_m \pi_{mk}\,x_{mk}^s\,\ln x_{mk}
    \;-\; \frac{B_k\,\ln B_k}{s}
\end{equation}
For the binary $L_1$, this reduces to:
\begin{equation}
    R_1 = (1 - w_{0.50})\cdot 10^s\,\ln 10
    + w_{0.50}\cdot 50^s\,\ln 50
    - \frac{B_1\,\ln B_1}{s}
\end{equation}

\subsection{The 7 Principal Minors}

Denote the entries of $\widetilde{J}^{\,\mathrm{SA}}$ by
$\widetilde{J}^{\,\mathrm{SA}}_{k\ell}$ where $k$ indexes the lottery
(row) and $\ell$ indexes the parameter (column: $r$, $\phi$, or $\eta$).
The P-matrix condition on $\widetilde{J}^{\,\mathrm{SA}}$ requires all 7
principal minors to be strictly positive:
\begin{enumerate}
    \item \textbf{Order 1 (3 conditions):} The diagonal entries, each
    representing the response of a lottery's CE to its assigned parameter
    (after sign correction):
    \begin{align}
        M_1 &= \widetilde{J}^{\,\mathrm{SA}}_{1,r} > 0
        &\quad&\text{($L_1$ responds to $r$)} \\
        M_2 &= \widetilde{J}^{\,\mathrm{SA}}_{2,\phi} > 0
        &\quad&\text{($L_2$ responds to $\phi$)} \\
        M_3 &= \widetilde{J}^{\,\mathrm{SA}}_{3,\eta} > 0
        &\quad&\text{($L_3$ responds to $\eta$)}
    \end{align}

    \item \textbf{Order 2 (3 conditions):} Each $2 \times 2$ principal
    minor must be positive, meaning pairwise own-effects dominate
    cross-effects:
    \begin{align}
        M_{12} &= \widetilde{J}^{\,\mathrm{SA}}_{1,r}\;
                   \widetilde{J}^{\,\mathrm{SA}}_{2,\phi}
                 - \widetilde{J}^{\,\mathrm{SA}}_{1,\phi}\;
                   \widetilde{J}^{\,\mathrm{SA}}_{2,r} > 0 \\
        M_{13} &= \widetilde{J}^{\,\mathrm{SA}}_{1,r}\;
                   \widetilde{J}^{\,\mathrm{SA}}_{3,\eta}
                 - \widetilde{J}^{\,\mathrm{SA}}_{1,\eta}\;
                   \widetilde{J}^{\,\mathrm{SA}}_{3,r} > 0 \\
        M_{23} &= \widetilde{J}^{\,\mathrm{SA}}_{2,\phi}\;
                   \widetilde{J}^{\,\mathrm{SA}}_{3,\eta}
                 - \widetilde{J}^{\,\mathrm{SA}}_{2,\eta}\;
                   \widetilde{J}^{\,\mathrm{SA}}_{3,\phi} > 0
    \end{align}

    \item \textbf{Order 3 (1 condition):}
    \begin{equation}
        M_{123} = \det \widetilde{J}^{\,\mathrm{SA}} > 0
    \end{equation}
\end{enumerate}

All entries of $\widetilde{J}^{\,\mathrm{SA}}$ are explicit functions of
$(r, \phi, \eta)$ through the formulas of the preceding subsection:
$B_k$ and its derivatives $\partial B_k/\partial\phi$,
$\partial B_k/\partial\eta$ involve only elementary operations
(powers $x^s$, exponentials, logarithms) applied to the fixed prize
and probability data of each lottery. The $r$-column involves the
additional quantity $R_k$ from~\eqref{eq:r_col}. There are no implicit
equations to solve---unlike the general $\widetilde{J}$ of
Section~\ref{sec:jacobian}, which requires computing CEs implicitly,
the power-mean representation makes everything explicit.

\subsection{Shared Prizes: P-Matrix Failure}

When all three lotteries share prizes $(a, b, c)$, the companion paper's
numerical analysis (the \texttt{injectivity\_analysis.py} code) has
demonstrated:
\begin{itemize}
    \item $\det \widetilde{J} \neq 0$ at generic parameter values (local
    injectivity).
    \item The mapping is \emph{not} globally injective: distinct $(r, \phi,
    \eta)$ triples produce identical CE vectors.
\end{itemize}

This is consistent with a failure of the P-matrix condition. Specifically, the
$2 \times 2$ minor $M_{23}$ involving the two Prelec columns is vulnerable:
with shared prizes, the $B$-block's rows differ only through the decision
weights, and the Prelec derivatives
$\partial w / \partial \phi$ and $\partial w / \partial \eta$ create a
specific interaction pattern that allows $M_{23}$ to change sign across the
parameter space. Where $M_{23} \leq 0$, the P-matrix condition fails, and
the Gale--Nikaido theorem cannot certify injectivity. The numerical finding
of non-injectivity confirms that this failure is genuine, not an artifact.

\subsection{Prize Diversity: P-Matrix Restoration}

The same numerical analysis demonstrates that using lotteries with
\emph{different prize vectors} restores global injectivity. In the framework
of this note, the mechanism is clear: different prizes give each row of
$\widetilde{J}$ a distinct dependence on the utility function and its
derivatives, which generically makes all principal minors positive.

The companion code's ``lottery design analysis'' (Analysis 6 in
\texttt{injectivity\_analysis.py}) shows that:
\begin{itemize}
    \item Adding a 4th lottery on the same prizes ($K > p$: overdetermination)
    also restores injectivity, because the system is overconstrained.
    \item Using even one lottery with different prizes suffices to break the
    compensation mechanism.
\end{itemize}

Both findings are consistent with the P-matrix framework: overdetermination
provides more constraints than degrees of freedom, while prize diversity
differentiates the Jacobian rows.


%==============================================================================
\section{Limitations}
\label{sec:limitations}
%==============================================================================

\paragraph{Rectangular domain requirement.}
The Gale--Nikaido theorem requires a rectangular (product of intervals) domain.
Economic parameter spaces are typically of this form---each parameter has its
own range, independent of others---so this is rarely a binding constraint.
However, if the model imposes \emph{joint} constraints on parameters (e.g.,
$\theta_1 + \theta_2 < 1$), the domain is not rectangular and the theorem
does not apply directly. In such cases, one may attempt to reparameterize to
obtain a rectangular domain, or use alternative injectivity results (Hadamard,
topological degree theory).

\paragraph{Unbounded domains.}
The theorem permits unbounded intervals ($a_i = -\infty$ or $b_i = +\infty$).
However, verifying the P-matrix condition on an unbounded domain requires
establishing that all principal minors remain positive as parameters approach
their extremes. In practice, one truncates to a compact rectangle covering the
prior support (e.g., a high-probability region under the prior) and verifies the
P-matrix condition there. This gives injectivity on the truncated domain, which
is sufficient for the Jacobian correction methodology since the prior assigns
negligible mass outside the truncated region.

\paragraph{Not a necessary condition.}
The P-matrix condition is sufficient but not necessary for global injectivity.
A mapping can be globally injective on a rectangular domain even if its
Jacobian is not a P-matrix at every point. The theorem provides a one-way
implication: P-matrix everywhere $\Rightarrow$ injective. The converse is
false. Thus, failure to verify the P-matrix condition does not prove
non-injectivity---it merely means this particular tool is inconclusive.

\paragraph{Computational cost of verification.}
The number of principal minors grows as $2^p - 1$, which is exponential in
the number of parameters. For small $p$ (the companion paper's case with
$p = 3$ has only 7 minors), this is trivial. For larger models (e.g., $p = 10$
would require 1023 minors), direct verification becomes expensive. The
diagonal dominance sufficient condition (\Cref{prop:diag_dom}) provides a
polynomial-cost alternative, at the price of a stronger requirement.

\paragraph{Dependence on lottery design.}
The P-matrix property depends on the choice of reference lotteries. A poor
choice of lotteries (e.g., all with identical prizes and similar probability
structures) can cause the P-matrix condition to fail even when a better design
would satisfy it. The researcher must therefore \emph{design} the reference
lotteries with injectivity in mind---this is a feature, not a bug, as it
provides actionable guidance for experimental design.


%==============================================================================
\section{Connection to the Paper's Appendix C}
\label{sec:connection}
%==============================================================================

The companion paper's Appendix C currently states (paraphrasing lines
1515--1519):

\begin{quote}
The classical sufficient conditions---Hadamard's theorem (properness) and the
Gale--Nikaido theorem (P-matrix Jacobian)---do not directly apply because CE
mappings are never proper (the codomain is bounded by prize values while the
parameter domain is unbounded).
\end{quote}

This statement conflates two distinct requirements:
\begin{itemize}
    \item \textbf{Hadamard's theorem} requires properness (or coercivity):
    $\|g(\bxi)\| \to \infty$ as $\|\bxi\| \to \infty$. CE mappings indeed
    fail this condition, since CEs are bounded by the prize range. Thus
    Hadamard's theorem genuinely does not apply.

    \item \textbf{The Gale--Nikaido theorem} requires a P-matrix Jacobian on
    a rectangular domain. It does \emph{not} require properness. The domain
    requirement is naturally satisfied by economic parameter spaces, and the
    P-matrix condition is a verifiable algebraic property of the Jacobian.
\end{itemize}

The corrected assessment is:
\begin{keyresult}[Corrected assessment of classical theorems]
\begin{itemize}
    \item \textbf{Hadamard:} Does not apply (CE mappings are bounded, so
    properness fails).
    \item \textbf{Gale--Nikaido:} \emph{Can} apply. The theorem's domain
    requirement (rectangular region) is satisfied by construction. The
    operative condition is the P-matrix property of the Jacobian, which is
    verifiable and is generically satisfied when reference lotteries have
    sufficient prize diversity.
\end{itemize}
\end{keyresult}


%==============================================================================
\section{Summary}
\label{sec:summary}
%==============================================================================

\begin{keyresult}[Main result]
The Gale--Nikaido univalence theorem provides the formal link between
\emph{prize diversity} (a lottery design choice controlled by the researcher)
and \emph{global injectivity} (the mathematical requirement for the Jacobian
correction to produce a valid probability density).

The chain of reasoning is:
\begin{enumerate}
    \item The CE mapping $g: \Xi \to \R^K$ maps parameters to certainty
    equivalents.
    \item The Jacobian $J_g$ factors as $D^{-1} \widetilde{J}$ with $D$
    positive diagonal.
    \item $J_g$ is a P-matrix iff $\widetilde{J}$ is a P-matrix.
    \item Prize diversity differentiates the rows of $\widetilde{J}$, generically
    ensuring all $2^p - 1$ principal minors are positive.
    \item By the Gale--Nikaido theorem, the P-matrix property on the rectangular
    parameter domain $\Xi$ implies global injectivity of $g$.
    \item Global injectivity validates the Jacobian correction formula
    $p_{\bxi}(\bxi) = p_{\mathbf{Y}}(g(\bxi)) \, |\det J_g(\bxi)|$.
\end{enumerate}
\end{keyresult}

\noindent
Key points:
\begin{itemize}
    \item The Gale--Nikaido theorem \emph{does} apply to CE mappings, contrary
    to the paper's current statement. The theorem's requirement is a rectangular
    domain and a P-matrix Jacobian---not properness (which is Hadamard's
    requirement).

    \item The P-matrix condition is the operative constraint. It is verifiable
    analytically (for small $p$) or numerically (via grid evaluation or interval
    arithmetic).

    \item Prize diversity is the economic mechanism that delivers the P-matrix
    property: distinct prize vectors differentiate the Jacobian's rows, creating
    the approximate diagonal dominance that implies the P-matrix condition.

    \item For the paper's 3-parameter case ($r, \phi, \eta$), only 7 principal
    minors need to be checked. The companion code's numerical analysis confirms
    that shared prizes cause P-matrix failure (and non-injectivity), while prize
    diversity restores it.

    \item The theorem provides \emph{sufficient} conditions. Failure to verify
    the P-matrix condition does not prove non-injectivity, but the companion
    paper's numerical analysis provides independent evidence of non-injectivity
    when the conditions fail.
\end{itemize}


%==============================================================================
% References
%==============================================================================

\begin{thebibliography}{9}

\bibitem{gale_nikaido_1965}
Gale, D.\ and Nikaido, H.\ (1965).
``The Jacobian matrix and global univalence of mappings.''
\textit{Mathematische Annalen}, 159(2), 81--93.

\bibitem{horn_johnson_2012}
Horn, R.\,A.\ and Johnson, C.\,R.\ (2012).
\textit{Matrix Analysis}, 2nd ed.
Cambridge University Press.

\bibitem{fiedler_ptak_1962}
Fiedler, M.\ and Pt\'ak, V.\ (1962).
``On matrices with non-positive off-diagonal elements and positive
principal minors.''
\textit{Czechoslovak Mathematical Journal}, 12(3), 382--400.

\bibitem{parthasarathy_1983}
Parthasarathy, T.\ (1983).
\textit{On Global Univalence Theorems}.
Lecture Notes in Mathematics, Vol.\ 977.
Springer-Verlag.

\bibitem{mas_colell_1979}
Mas-Colell, A.\ (1979).
``Homeomorphisms of compact, convex sets and the Jacobian matrix.''
\textit{SIAM Journal on Mathematical Analysis}, 10(6), 1105--1109.

\end{thebibliography}

\end{document}
